% Self-contained proof of cp:initialization, stated in compact.tex.
\section{Efficient initialization of the selected models}
\label{cp:init:proof}

We prove \Cref{cp:initialization} by computing disposable source values
on $[0,T]$, $T=32(m/\gamma)\log(en)$, then applying \Cref{cp:selected}
to real-node interpolation (Taylor) or positive coefficient quadrature
(Harmonic). The autonomous model starts at time zero with original
activations, exact initialized training features and
\eqref{cp:selected-readout}. Supplied matrices specify its dense reference;
the implicit algorithm instead samples a fresh Gaussian reference through
consistent forward and transpose actions. It does not convert supplied
matrices or a prescribed entrywise seed in subquadratic work.

\subsection{Numerical source values from initialization}
\label{cp:init:sec:numerics}

All constants here are local to the numerical proof. Put $\lambda=\gamma/m$, $S=16Y/\lambda$,
$\ell=\log(en)$, $T=32\ell/\lambda$, and $W_0^{(j)}=W^{(j)}(0)$.
The normalized parameter vector and its norm are

\[
\theta=(W^{(1)}/\sqrt n,W^{(2)},\ldots,W^{(L)},w/\sqrt n),
\qquad
\|\theta\|^2=\|W^{(1)}\|_F^2/n+
 \sum_{j=2}^L\|W^{(j)}\|_F^2+\|w\|_2^2/n.
\]

For the predictions and residuals in \eqref{cp:network}, set
$\xi_a=\nabla_\theta f_n(x_a)$ and
$\rho=\|r\|_2/\sqrt m$. The mobilities in
\eqref{cp:flow} give precisely

\begin{equation}
\label{cp:init:eq:field}
\dot\theta=F(\theta),\qquad
F(\theta)=-\frac2m\sum_a r_a\xi_a,\qquad
\xi_a=\left(\frac{\delta_a^{(1)}v_a^\top}{\sqrt n},
 \left(\frac{\delta_a^{(j)}h_a^{(j-1)\top}}n\right)_{j=2}^L,
 \frac{h_a^{(L)}}{\sqrt n}\right).
\end{equation}

Complex products use algebraic transposes and ordinary complex
Euclidean/Frobenius norms; no complex Gram positivity is used.

Condition on \Cref{cp:fit,cp:source} for Harmonic and
\Cref{cp:poly:source} for Taylor. These events give real operator caps below nine,
$\|w\|_2/\sqrt n\le2Y/\sqrt\lambda$,
$\rho(t)\le Ye^{-\lambda t/2}$ and $\int_0^\infty\rho\le2Y/\lambda$.
For centers in $[0,T]$, the respective source proofs supply complex
disks of radius $r_t/2$ with preactivation imaginary parts at most
$3a/8$, operator caps ten, residual RMS at most $2Y$, feature RMS at most
$\beta^{3L}$, and training carrier coordinates bounded by

\begin{equation}
\label{cp:init:eq:source-input}
r_t=\frac{\lambda}{\beta^{30L}Y^2\sqrt{(d+3)\ell}},\qquad
M_c=2S\beta^{21L}\sqrt\ell.
\end{equation}

Here $k^{(L)}=w$, $k^{(j)}=W^{(j+1)\top}\delta^{(j+1)}$.
The displayed $r_t$ is for Harmonic. Taylor uses $r_t=h(0)/2$
from \eqref{cp:init:eq:capped-panels}; its radius-$h(0)/4$ disk
is certified at every real center. Only declared Taylor queries, or real
unit Harmonic queries, use these complex domains. For Taylor, apply the
input reduction below and use its reduced dimension in the preactivation
envelope.

\subsubsection{Real signed stability and its constants}
\label{cp:init:sec:stability}

Use $b=\beta-1$, $s=t_2=\beta$, and define in order

\begin{equation}
\label{cp:init:eq:real-constants}
\begin{gathered}
H_1^{\rm r}=\max(1,b+10s),\quad
H_j^{\rm r}=\max(1,b+10sH_{j-1}^{\rm r}),\quad H^{\rm r}=H_L^{\rm r},\quad
R^{\rm r}=1+2Y/\sqrt\lambda,\\
B^{\rm r}=R^{\rm r}s(10s)^{L-1},\quad
F_z^{\rm r}=H^{\rm r}(10s)^{L-1},\quad
P_{\rm layer}^{\rm r}=1+(L-1)H^{\rm r},\\
D^{\rm r}(M)=(10s)^{L-1}
 [s(1+B^{\rm r})+Lt_2F_z^{\rm r}M],\\
J^{\rm r}(M)=\sqrt{L+1}
 [P_{\rm layer}^{\rm r}D^{\rm r}(M)+(1+(L-1)B^{\rm r})sF_z^{\rm r}],\\
C^{\rm r}(M)=sF_z^{\rm r}\sqrt L+LB^{\rm r}sF_z^{\rm r}
                  +\tfrac12L^2t_2(F_z^{\rm r})^2M.
\end{gathered}
\end{equation}

Evaluate $J^{\rm r},C^{\rm r}$ at $M=M_c$, and put

\begin{equation}
\label{cp:init:eq:real-source-constant}
\begin{gathered}
A_n=4J^{\rm r}Y/\lambda,\quad
K_{\rm qry}^{\rm r}=R^{\rm r}(10s)^{L-1},\\
C_{\rm src}=8\sqrt{L+1}\max\left\{sF_z^{\rm r},
 (10s)^{L-1}[s(1+B^{\rm r})+Lt_2F_z^{\rm r}K_{\rm qry}^{\rm r}]\right\}.
\end{gathered}
\end{equation}

\begin{lemma}[Stability of an approximate dense path]
\label{cp:init:defect}

Let $u$ be an absolutely continuous path on $[0,T]$, let
$d_u=u'-F(u)$, and set
$E_0=\|u(0)-\theta(0)\|+\int_0^T\|d_u\|$.
If

\[
E_0\le e^{-A_n}/4,\qquad
8(C^{\rm r}+J^{\rm r})^2Te^{2A_n}E_0^2\le1,
\]

then $\sup_{[0,T]}\|u-\theta\|\le2e^{A_n}E_0$.
Every real-query source in \eqref{cp:source-families} has coordinate
error at most $C_{\rm src}n\|u(t)-\theta(t)\|$ when evaluated at
$u(t)$ with the original activations.
\end{lemma}

\begin{proof}
Stop at $E=\|u-\theta\|=1$, so both states have operator cap ten
and readout RMS cap $R^{\rm r}$. For the sum $D_h$ of hidden-block
discrepancies and readout RMS discrepancy $D_w$,
$D_h\le\sqrt L E$, $D_h+D_w\le\sqrt{L+1}E$.
Forward subtraction bounds preactivation and feature RMS errors by
$F_z^{\rm r}D_h$ and $sF_z^{\rm r}D_h$. Backward subtraction is

\[
\delta^{(j)}(u)-\delta^{(j)}(\theta)
=\phi'(z^{(j)}(u))\odot[k^{(j)}(u)-k^{(j)}(\theta)]
 +[\phi'(z^{(j)}(u))-\phi'(z^{(j)}(\theta))]\odot k^{(j)}(\theta).
\]

The second term uses $M_c$; propagation, changed mixers and top readout
cost $10s$, $sB^{\rm r}\|\Delta W\|_F$ and $sD_w$.
Their sum gives response RMS error $D^{\rm r}(M_c)(D_h+D_w)$
and $\|\xi_a(u)-\xi_a(\theta)\|\le J^{\rm r}E$.
The endpoint remainder obeys

\[
|f_n(u,x_a)-f_n(\theta,x_a)-\langle\xi_a(\theta),u-\theta\rangle|
 \le C^{\rm r}E^2.
\]

Indeed, put
$a^{(j)}=\Delta h^{(j)}-\phi'(z^{(j)}(\theta))\odot\Delta z^{(j)}$.
Scalar Taylor's formula gives
$|a_i^{(j)}|\le t_2|\Delta z_i^{(j)}|^2/2$.
Backward substitution expresses the remainder as

\[
\frac1n\sum_j k_a^{(j)}(\theta)^\top a_a^{(j)}
 +\frac1n\sum_{j=2}^L\delta_a^{(j)}(\theta)^\top
                  \Delta W^{(j)}\Delta h_a^{(j-1)}
 +\frac1n\Delta w^\top\Delta h_a^{(L)}.
\]

Its three bounds are
$L^2t_2M_c(F_z^{\rm r})^2E^2/2$,
$LB^{\rm r}sF_z^{\rm r}E^2$, and $sF_z^{\rm r}\sqrt L E^2$.
For $e_{f,a}=f_n(u,x_a)-f_n(\theta,x_a)$, the exact signed identity
in the proof of \Cref{cp:signed}, with these two remainder bounds, gives

\[
\tfrac12(E^2)'
\le-2\|e_f\|_2^2/m
 +2(C^{\rm r}+J^{\rm r})\|e_f\|_2E^2/\sqrt m
 +2J^{\rm r}\rho E^2+E\|d_u\|.
\]

Maximizing over $\|e_f\|_2/\sqrt m$ gives

\begin{equation}
\label{cp:init:eq:signed-defect}
E'\le2J^{\rm r}\rho E+\tfrac12(C^{\rm r}+J^{\rm r})^2E^3+\|d_u\|.
\end{equation}

The regularization in \Cref{cp:signed} covers $E=0$.
Multiplying by $\exp(-2J^{\rm r}\int_0^t\rho)$, whose inverse is
at most $e^{A_n}$, bounds the transformed error before its first
$2E_0$ crossing by
$E_0+4(C^{\rm r}+J^{\rm r})^2Te^{2A_n}E_0^3\le3E_0/2$.
This excludes the crossing and gives $E\le1/2$, removing the stop;
the same inequality covers $E_0=0$.
For any real query, replace the carrier cap by
$\sqrt nK_{\rm qry}^{\rm r}$ in the same subtraction and convert
RMS to coordinates. These two $\sqrt n$ factors and initialized
operator cap eight give $C_{\rm src}nE$, also between any two
states in this real neighborhood. Finally $A_n=O(\sqrt\ell)$
at fixed parameters.
\end{proof}

\subsubsection{Complex restarts and scalar activation approximation}
\label{cp:init:sec:activation}

For complex neighborhoods of computed anchors, define separately

\begin{equation}
\label{cp:init:eq:complex-constants}
\begin{gathered}
R_{\rm loc}=1+S\beta^{3L},\quad
H_1=\max(1,b+11s),\quad H_j=\max(1,b+11sH_{j-1}),\\
B_j=R_{\rm loc}s(11s)^{L-j},\quad
P_1=1,\quad P_j=H_{j-1}+11sP_{j-1},\quad P_*=\max_jP_j,\\
\kappa_L(M)=1,\quad
\kappa_j(M)=11s\kappa_{j+1}(M)+11t_2MP_{j+1}+B_{j+1},\\
D_j(M)=s\kappa_j(M)+t_2MP_j,\quad \kappa_*(M)=\max_j\kappa_j(M),\\
G=H_L+B_1+\sum_{j=2}^LB_jH_{j-1},\\
J_{\rm loc}(M)=sP_L+D_1(M)
  +\sum_{j=2}^L[H_{j-1}D_j(M)+B_jsP_{j-1}].
\end{gathered}
\end{equation}

The same endpoint subtraction, now at operator cap eleven and readout
RMS cap $R_{\rm loc}$, gives preactivation, carrier and response RMS
errors $P_jE,\kappa_j(M)E,D_j(M)E$ when the reference carrier
maximum is $M$. The gradient blocks then give $\|\xi_a\|\le G$,
$\|\Delta\xi_a\|\le J_{\rm loc}(M)E$, and

\[
\|F(u)-F(v)\|\le[2G^2+2\rho(v)J_{\rm loc}(M)]\|u-v\|.
\]

Here integration along the parameter segment bounds the residual change
by $GE$; all scalar preactivation segments lie in the half-strip below.

Define

\begin{equation}
\label{cp:init:eq:restart-radius}
\begin{gathered}
b_n=\min\{1/4,a/(16\sqrt nP_*),1/(\sqrt n\kappa_*(M_c))\},\\
L_n=2G^2+2(2Y+G)J_{\rm loc}(M_c+1),\quad V=2G(2Y+G),\quad
\widehat R_n=\min\{r_t/2,(8L_n)^{-1}\}.
\end{gathered}
\end{equation}

In a radius-$b_n$ ball about the true complex trajectory, first-exit
subtraction gives preactivation change at most $a/16$, hence imaginary
part below $7a/16$, operator cap eleven, carrier maximum $M_c+1$
and residual RMS $2Y+G$. Thus $L_n$ is a proved pairwise Lipschitz
bound there. The source proofs also give the center envelope

\[
Z_*=[\beta^{21L}+16\sqrt{d+3}\,\beta^{3L}+1
                 +\beta^{26L}\sqrt{d+3}\,S^2]\sqrt\ell+a/8.
\]

by integrating $2\rho S\beta^{26L}\sqrt{d+3}\sqrt\ell$
and adding the initial Gaussian and short complex-time bounds.
For Taylor use its finite-list maximum and
$SU_{\rm fin}\sqrt\ell$, $U_{\rm fin}\le\beta^{26L}$.
The moving-ball rectangle is therefore
$|\Re z|\le Z_*+a/16$, $|\Im z|\le7a/16$.
For source-vector tolerance $\varepsilon_v>0$, put
$\tau=\varepsilon_v/(32\sqrt n)$ and define

\begin{equation}
\label{cp:init:eq:activation-tolerance}
\begin{gathered}
e_d=e^{-A_n}\min\left\{\frac14,
 \frac1{\sqrt8(C^{\rm r}+J^{\rm r})\sqrt T},
 \frac{\tau}{4C_{\rm src}n},\frac{b_n}{16}\right\},\quad
\zeta_F=\min\{e_d/(2T),b_nL_n/4,V/2\},\\
U_0=0,\quad U_j=1+11sU_{j-1},\quad U_*=\max_jU_j,
 \quad K_j^{\rm act}=R_{\rm loc}(11s)^{L-j},\\
Q_{L+1}(M)=0,\quad
Q_j(M)=11(s+1)Q_{j+1}(M)+K_j^{\rm act}+11t_2MU_{j-1},\\
Q_\xi=U_L+Q_1(M_c+1)
 +\sum_{j=2}^L[Q_j(M_c+1)(H_{j-1}+1)+B_jU_{j-1}],\\
Q_f=R_{\rm loc}U_L,\quad
C_F=2Q_f(G+1)+2(2Y+G)Q_\xi,\\
M_q=\sqrt n\max_jK_j^{\rm act},\quad
C_p=8\sqrt n\max\{U_*,\max_jQ_j(M_q)\},\\
\epsilon=\min\{1,U_*^{-1},a/(5632\sqrt nU_*),Q_\xi^{-1},
                         \zeta_F/C_F,\tau/(2C_p)\}.
\end{gathered}
\end{equation}

Every denominator is positive. Set

\begin{equation}
\label{cp:init:eq:activation-degree}
\begin{gathered}
B=8(Z_*+a+1),\quad
\tau_o=\operatorname{arsinh}(31a/(64B)),\quad
\tau_i=\operatorname{arsinh}(29a/(64B)),\quad\Delta=\tau_o-\tau_i,\\
M_\phi=b+s(B+a),\quad C_a=\max\{1,512/a,2(512/a)^2\},
 \quad\epsilon_p=\epsilon/C_a,\\
D=\max\left\{1,\left\lceil\Delta^{-1}
       \log\max\{e,24M_\phi/(\epsilon_p\Delta)\}\right\rceil\right\},
 \qquad N_\phi=4(D+1).
\end{gathered}
\end{equation}

With $x_i=B\cos(2\pi i/N_\phi)$, compute for each layer

\begin{equation}
\label{cp:init:eq:activation-polynomial}
\begin{gathered}
c_{j,0}=N_\phi^{-1}\sum_i\phi_j(x_i),\qquad
c_{j,k}=2N_\phi^{-1}\sum_i\phi_j(x_i)\cos(2\pi ki/N_\phi) (k>0),\\
\psi_j(z)=\sum_{k=0}^Dc_{j,k}T_k(z/B).
\end{gathered}
\end{equation}

Finite value-call errors of magnitude at most
$\epsilon_p/[4(D+1)e^{D\tau_i}]$ are allowed in this formula only.
Initialized training features below still use exact original values.

\begin{lemma}[Certified numerical source polynomials]
\label{cp:init:source-routine}

Let $C_g=\max_j\{H_j+1,B_j+1\}$ and choose

\begin{equation}
\label{cp:init:eq:numerical-orders}
\begin{split}
K=\max\bigg\{&8,
 \left\lceil2\log_2\max\{1,16TV/e_d\}\right\rceil,
 \left\lceil\log_2\max\{1,8V\widehat R_n/b_n\}\right\rceil,\\
&\left\lceil\log_2\max\{1,128V\widehat R_nC_{\rm src}n^{3/2}/\varepsilon_v\}\right\rceil,
 \left\lceil\log_2\max\{1,64\sqrt nC_g/\varepsilon_v\}\right\rceil\bigg\},\\
J&=\lceil2T/\widehat R_n\rceil
 =\lceil\max\{4T/r_t,16TL_n\}\rceil.
\end{split}
\end{equation}

Use $J$ equal physical-time panels. Starting at the original initialization,
compute Taylor coefficients through degree $K$ for the polynomial-activation
field $\widetilde F$ by

\begin{equation}
\label{cp:init:eq:jet-recursion}
[u]_{k+1}=\frac1{k+1}[h^k]\widetilde F
                   \left(\sum_{i=0}^k[u]_ih^i\right),\qquad 0\le k<K,
\end{equation}

and advance to each polynomial endpoint. The forward and backward source
jets through degree $K$, computed with the same coefficients and
$\psi_j$, approximate every base source $h^{(j)}$ or
$\delta^{(j)}$ in Euclidean norm by at most $\varepsilon_v/16$
at every real time in its numerical panel. Queries are all real unit
inputs for Harmonic and the declared input list for Taylor.
If $\log(1/\varepsilon_v)=O(\ell)$ at fixed parameters, then

\begin{equation}
\label{cp:init:eq:numerical-scales}
J=O(\ell^{3/2}),\qquad K=O(\ell),\qquad D=O(\ell^{3/2}).
\end{equation}
\end{lemma}

\begin{proof}
The outer ellipse $z=B(\zeta+\zeta^{-1})/2$,
$|\zeta|=e^{\tau_o}$ lies in the half-strip and has
$|\phi|\le M_\phi$. The inner ellipse contains radius-$a/512$
disks about the $a/512$ neighborhood of the moving-ball rectangle,
since $(1/8)^2+(226/232)^2<1$; also
$a/(64B)\le\Delta<1$. The Fourier contour argument of
\Cref{cp:harmonic-approximation}, applied to $\phi(B\cos u)$,
gives coefficients $2M_\phi e^{-k\tau_o}$ and inner-ellipse tail
$4M_\phi\Delta^{-1}e^{-(D+1)\Delta}$.
Discrete sampling aliases frequencies $k+qN_\phi$; their geometric
sum for $k\le D$ is at most
$8M_\phi(D+1)e^{-2(D+1)\tau_o}
\le8M_\phi\Delta^{-1}e^{-(D+1)\Delta}$.
The prescribed $D$ makes these errors at most $\epsilon_p/2$;
summing the allowed value-call errors adds at most $\epsilon_p/2$.
Cauchy on the verified disks gives value and first-two-derivative
errors at most $\epsilon$.

At a fixed moving-ball state, forward substitution errors are
$U_j\epsilon$ in feature RMS and $11U_{j-1}\epsilon$ in
preactivation RMS. Since $11\sqrt nU_*\epsilon\le a/512$,
all polynomial arguments remain in that domain. Backward subtraction
has the three terms $11(s+1)Q_{j+1}$, $K_j^{\rm act}$ and
$11t_2MU_{j-1}$, proving the response bound $Q_j(M)\epsilon$.
The prediction and gradient bounds are $Q_f\epsilon,Q_\xi\epsilon$;
product subtraction in \eqref{cp:init:eq:field} gives
$\|\widetilde F-F\|\le C_F\epsilon\le\zeta_F$.
For passive sources and initialized images the coordinate error is
$C_p\epsilon\le\tau/2$. On the concentric half-ball,
directional Cauchy gives derivative error at most
$2\zeta_F/b_n\le L_n/2$, so $\widetilde F$ is $2L_n$-Lipschitz.

For $\|u_0-\theta(t)\|\le b_n/8$, consider on
$|h|\le\widehat R_n$ the holomorphic-error map

\[
e(h)\longmapsto u_0-\theta(t)+
 \int_0^h[\widetilde F(\theta(t+z)+e(z))-F(\theta(t+z))]\,dz.
\]

On the supremum ball of radius $b_n/2$ it contracts by at most
$1/4$ and has image norm at most $9b_n/32$.
Uniform Picard convergence and Cauchy's formula give a holomorphic
solution $v$, real by conjugation and uniqueness, with discrepancy
at most $5b_n/24<b_n/4$ and velocity at most $2V$.
Its degree-$K$ truncation $u$ on a half-radius panel has parameter
and derivative tails $2V\widehat R_n2^{-K}$ and
$2V(2K+4)2^{-K}$. It remains in the smaller moving ball, with
defect at most $2V(2K+5)2^{-K}+\zeta_F$.
Because $2K+5\le4\,2^{K/2}$ for $K\ge8$, the prescribed order
bounds total integrated defect by $e_d$. \Cref{cp:init:defect}
on every completed prefix gives

\[
\|u-\theta\|\le2e^{A_n}e_d
 \le\min\{b_n/8,\tau/(2C_{\rm src}n)\}.
\]

This closes finite restart induction from exact initialization, without
unknown-norm tests or rejected panels. For a base source let $g^0,g^p$
denote original and polynomial activation evaluation and $g_K$ the
degree-$K$ source polynomial along $v$. Its error splits as

\[
g_K-g^0(\theta)=[g_K-g^p(v)]+[g^p(v)-g^0(v)]
              +[g^0(v)-g^0(u)]+[g^0(u)-g^0(\theta)].
\]

Each term is at most $\varepsilon_v/64$: the first uses
$\|g^p(v)\|_2\le\sqrt n C_g$ and its geometric Taylor tail;
the second and fourth use $\sqrt n\tau/2$; the third uses
$2V\widehat R_nC_{\rm src}n^{3/2}2^{-K}$.
Here the bound $C_g$ follows from $H_j,B_j$ and substitution
error below one. This proves the asserted $\varepsilon_v/16$.
Finally $M_c,L_n,A_n=O(\sqrt\ell)$,
$b_n^{-1}=O(\sqrt{n\ell})$, $\widehat R_n^{-1}=O(\sqrt\ell)$,
$C_p=O(n)$ and $Z_*=O(\sqrt\ell)$ at fixed parameters.
Thus all inverse-tolerance logarithms are $O(\ell)$ and
$\Delta^{-1}=O(\sqrt\ell)$, proving \eqref{cp:init:eq:numerical-scales}.
The causal recursion differentiates only the finite polynomials $\psi_j$.
\end{proof}

\subsection{Two source spaces with the original retained sizes}
\label{cp:init:sec:source-spaces}

Let $\mathcal A_n$ be the comparison coefficient in
\eqref{cp:selected-coefficients}, and let $\mathcal D$ be the
explicit endpoint coefficient in the proof of
\Cref{cp:selected-fitting}. Use

\begin{equation}
\label{cp:init:eq:transfer-tolerance}
\eta=\min\{1,Y,S,Y/(2n\mathcal A_n)\},\qquad
\mathcal D e^{-8\ell}\le Y/(2n).
\end{equation}

Then \Cref{cp:selected} gives all-time error $Y/n$ once its source
contract holds. At fixed positive parameters, $\log(1/\eta)\le2\ell$
eventually. Taylor uses $C_{\rm sel}=\beta^{60L}$ from
\hyperref[cp:poly:assembly]{the polynomial-width assembly}: its conservative
tolerance is no larger than this $\eta$, and
\eqref{cp:poly:eq:taylor-tail-gate} gives the tail gate above.
We use exactly \eqref{cp:poly:eq:taylor-count-gate} for the degree,
so Taylor incurs no additional qualitative threshold. Harmonic retains
its eventual source and scalar gates.

\subsubsection{Taylor from real source nodes}
\label{cp:init:sec:taylor}

Put $M=m+p$ and apply the exact input reduction proved in \Cref{cp:panel}:
$U_0\in\R^{d\times k}$ spans the declared normalized inputs,
$U_0^\top U_0=I$, $k\le\min(d,M)$, and the reduced inputs and
first matrix are $U_0^\top v_i,W_0^{(1)}U_0$. This preserves
their Gram, Gaussian law and all declared trajectories; the retained
map costs at most $Md$.
Use the deterministic finite-panel domain of \Cref{cp:poly:source}, with
$\eta_*,U_{\rm fin},\mathcal K$ defined by
\hyperref[cp:src-coefficients]{the source recurrences}, \eqref{cp:src-UV}
and \eqref{cp:src-K}, without observed trajectory norms:

\begin{equation}
\label{cp:init:eq:capped-panels}
\begin{gathered}
\nu=\lambda/4,\quad a_0=\min(a,1),\quad
b_{\rm cap}=\log(e+\ell),\quad B_{\rm cap}=1+(S^2/\eta_*)b_{\rm cap},\\
r_{\rm cap}=\min\left\{\frac{a_0}{64YSU_{\rm fin}\sqrt\ell},
 \frac1{\beta^{300L}(1+\lambda B_{\rm cap})\sqrt{b_{\rm cap}}}\right\},\\
h(t)=\frac{\log(1+2\mathcal K r_{\rm cap}e^{\nu t})}{2\mathcal K},
\qquad d_h=\lambda/(8\mathcal K).
\end{gathered}
\end{equation}

Set $t_0=0$, $r_b=h(t_b)/[2(1+d_h)]$ and
$t_{b+1}=\min\{T,t_b+r_b/2\}$. These certified disks have radius
at least $4h(0)/9>h(0)/4$, since $h$ is nondecreasing and
$d_h\le1/8$, verifying the numerical disk used above. Identifying
$X=\beta^L$, $\ell_n=\ell$, $r_n=r_{\rm cap}$ in
\eqref{cp:poly:eq:panel-count} gives the source-panel count

\begin{equation}
\label{cp:init:eq:taylor-panels}
J_s\le1+\beta^{400L}\left[(Y/\lambda)^2\sqrt\ell
                  +\lambda^{-1}\log(e+\ell)\right].
\end{equation}

Here also $b_n=b_{\rm cap}$ and $B_n=B_{\rm cap}$ in that proof;
its confidence-free count preserves the $(Y/\lambda)^2$ factor.
Since $2\mathcal K r_{\rm cap}\le1$, $h(0)\ge r_{\rm cap}/2$
and $h(0)^{-1}=O(\sqrt\ell)$ at fixed parameters. Thus
\eqref{cp:init:eq:numerical-scales} holds even while the cap is active.

Put $M_n=\beta^{6L}\sqrt n$ and

\begin{equation}
\label{cp:init:eq:taylor-nodes}
K_s=\max\{8,\lceil2\log_2(64M_n/\eta)\rceil\},\quad
N_s=J_s(K_s+1),\quad
\varepsilon_v=\eta/[128(K_s+1)].
\end{equation}

On each $[t_b,t_{b+1}]$, evaluate the $K_s+1$ Chebyshev roots
$t_b+(t_{b+1}-t_b)[1+\cos((i+1/2)\pi/(K_s+1))]/2$
by \Cref{cp:init:source-routine}, forward on all $M$ inputs and backward
on training inputs, and insert exact initialized images.
The nodes generate the space without computed interpolation coefficients.
Discrete cosine orthogonality and $|T_j|\le1$ give
$\|I_{K_s}\|\le2K_s+1$, in coordinate or Euclidean norm.
Reproduction of the degree-$K_s$ Taylor polynomial, whose tail on the
half-radius panel is $M_n2^{-K_s}$, gives
\[
\|g-I_{K_s}g\|_\infty
 \le(2K_s+3)M_n2^{-K_s}\le4M_n2^{-K_s/2}\le\eta/16
\]
for each exact source and image. Numerical nodes and images have
Euclidean errors at most $\varepsilon_v,8\varepsilon_v$, so their
interpolated perturbation is at most $16(K_s+1)\varepsilon_v=\eta/8$.
Both families therefore meet accuracy $\eta$, with exact pairing
$I_{K_s}(W_0\widetilde g)=W_0I_{K_s}\widetilde g$ and its transpose
counterpart.

Add the constant, exact initialized training features and images, and
first-layer columns. The rank and actual selected width obey
\begin{equation}
\label{cp:init:eq:taylor-rank}
r\le R=2m+k+1+2(M+m)N_s,\qquad q\le9r.
\end{equation}
Apply \Cref{cp:selected} with source-dimension bound $R$.
To verify this numerical degree at the stated polynomial gate, put
$B_0=32\beta^{6L}(1+C_{\rm sel})\lambda^{-3/2}(1+Y^{-1})$.
At $n\ge B_0^4$ and the source gate,
$\ell\ge16$, $\log B_0\le(\log n)/4$. On the active comparison
branch, $1+\lambda^{-1/2}\le2\lambda^{-1/2}$ gives
\[
\log(64M_n/\eta)
\le\tfrac32\log n+\log B_0+\log8+\sqrt\ell+\log(1+\sqrt\ell)
\le\tfrac94\ell+\log8.
\]
The $Y$ and $1$ branches are smaller; $S\ge16Y$ is inactive.
Thus $K_s+1\le\max\{9,(9/(2\log2))\ell+8\}\le8\ell$.
The rank squaring in \Cref{cp:panel}, now using
\eqref{cp:init:eq:taylor-panels}, proves \eqref{cp:panel-size}
with its $(Y/\lambda)^4\ell^3$ factor. The map and data add $O(Md)$;
no passive labels enter.
\subsubsection{Harmonic from positive coefficient quadrature}
\label{cp:init:sec:harmonic}

Use the current source radii $r_t$ in
\eqref{cp:init:eq:source-input} and
$r_q=1/[\beta^{3L}\sqrt{(d+3)\ell}]$, with
$M_n=\beta^{6L}\sqrt n$ and $\alpha=r_t/(4T)$.
Require the eventual scalar gates $\alpha,r_q\le1$.
For $d\ge2$, take exactly the mode set of
\Cref{cp:harmonic-approximation}:

\begin{equation}
\label{cp:init:eq:harmonic-modes}
\begin{gathered}
b_d=2d-2,\quad D_d=2^{d+1}d^{d-2},\quad
P_T=18D_db_d!2^{b_d+1}/(\alpha r_q^{b_d+1}),\quad
H_*=2\log(16M_nP_T/\eta),\\
\Lambda=\{(k,j,b):\alpha k+r_qj\le H_*,\ 1\le b\le h_j\},\quad
h_j={j+d-1\choose d-1}-{j+d-3\choose d-1},\\
N=|\Lambda|,\quad p_t=\lfloor H_*/\alpha\rfloor,\quad
j_*=\lfloor H_*/r_q\rfloor,\quad
H_{\rm sph}=\sum_{j\le j_*}h_j,\quad\mathcal Y=\max_{j\le j_*}\sqrt{h_j}.
\end{gathered}
\end{equation}

Impossible binomials are zero. For a scalar source $g$, its coefficient
is

\begin{equation}
\label{cp:init:eq:harmonic-coefficient}
c_{kjb}=\gamma_k\int_0^{2\pi}\!\int_{S^{d-1}}
g(T(1+\cos u)/2,v)\cos(ku)Y_{jb}(v)\,d\sigma(v)\frac{du}{2\pi},
\quad\gamma_0=1,\quad\gamma_k=2\ (k>0),
\end{equation}

where $Y_{jb}$ are real orthonormal harmonics for probability surface
measure. The proof of \Cref{cp:harmonic-approximation} bounds
the omitted-mode tail by $\eta/16$.

Here are finite quadrature orders and their construction. For $d\ge3$
use polar angles $\theta_1,\ldots,\theta_{d-2}$ and azimuth
$\varphi$ in the usual recursive coordinates
$v=(\cos\theta_1,\sin\theta_1u)$; for $d=2$ use only
$v=(\cos\varphi,\sin\varphi)$. Define

\begin{equation}
\label{cp:init:eq:quadrature-orders}
\begin{gathered}
I_0=\pi,\quad I_1=2,\quad I_b=(b-1)I_{b-2}/b,\quad
A_d=\prod_{b=1}^{d-2}\pi/I_b,\\
a_t=\alpha/2,\quad h_{\rm ang}=r_q/[2(d-1)],\quad
\sigma=\operatorname{arsinh}(2h_{\rm ang}/\pi),\\
\epsilon_c=\eta/(16N\mathcal Y),\quad
\delta_{\rm node}=\eta/(64NA_d\mathcal Y^2),\\
B_{\rm int}=2M_n\mathcal YA_d
 e^{a_tp_t+j_*(d-1)h_{\rm ang}}
                 (\cosh h_{\rm ang})^{(d-2)(d-1)/2},\\
N_t=\max\{1,\lceil a_t^{-1}\log(1+4dB_{\rm int}/\epsilon_c)\rceil\},\\
N_\varphi=\max\{1,\lceil h_{\rm ang}^{-1}\log(1+4dB_{\rm int}/\epsilon_c)\rceil\},\\
N_\theta=\max\{1,\lceil\sigma^{-1}
 \log[8dB_{\rm int}/(\epsilon_c(1-e^{-\sigma}))]\rceil\}\quad(d\ge3),\\
N_x=N_\varphi N_\theta^{d-2}\quad(d\ge3),\qquad N_x=N_\varphi\quad(d=2).
\end{gathered}
\end{equation}

Empty products equal one. Apply equal-weight periodic trapezoids with
$N_t,N_\varphi$ nodes to $u,\varphi$. For each polar coordinate,
transport to $\theta=\pi(1+x)/2$ the nodes and weights

\begin{equation}
\label{cp:init:eq:fejer}
\omega_i=(2i+1)\pi/(2N_\theta),\quad x_i=\cos\omega_i,\quad
w_i=\frac1{N_\theta}\left[1-2\sum_{k=1}^{\lfloor(N_\theta-1)/2\rfloor}
                    \frac{\cos(2k\omega_i)}{4k^2-1}\right].
\end{equation}

Multiply spatial weights by the normalized spherical Jacobian
$\prod_{i=1}^{d-2}(\pi/I_{d-1-i})\sin^{d-1-i}\theta_i$;
integration by parts verifies $I_b=\int_0^\pi\sin^b\theta\,d\theta$.
The Fourier estimate already used above gives trapezoid error
$2B/(e^{aN}-1)$ for strip half-width $a$ and bound $B$, since
the rule retains exactly the frequencies divisible by $N$.
Discrete cosine orthogonality gives the polar interpolant
$N_\theta^{-1}\sum_i g_i[1+2\sum_{k=1}^{N_\theta-1}
\cos(k\omega_i)T_k(x)]$.
Its normalized integral is \eqref{cp:init:eq:fejer}: odd modes
integrate to zero and $T_{2k}$ to $-1/(4k^2-1)$.
The rule is exact through $N_\theta-1$; its weights sum to one and
are positive because $2\sum_{k=1}^q(4k^2-1)^{-1}=1-(2q+1)^{-1}<1$.
Twice the ellipse Chebyshev tail therefore bounds its error by
$4Be^{-N_\theta\sigma}/(1-e^{-\sigma})$, since both positive rule
and normalized integral have norm one.

The domain checks hold jointly. A complex rotation of imaginary angle
$y$ sends imaginary norm $\sinh r$ to at most $\sinh(r+|y|)$:
the real and imaginary parts of a quadric point have norms
$\cosh r,\sinh r$, and the rotation costs $\cosh|y|,\sinh|y|$.
Thus the $d-1$ rotations stay in the radius-$r_q/2$ tube.
A harmonic has trigonometric degree at most $j_*$ in each angle,
so Laurent maximum modulus bounds it by
$\mathcal Ye^{j_*(d-1)h_{\rm ang}}$; the Jacobian contributes
$A_d(\cosh h_{\rm ang})^{(d-2)(d-1)/2}$, giving $B_{\rm int}$.
The polar ellipse has height $(\pi/2)\sinh\sigma=h_{\rm ang}$.
Telescoping the $d$ integral/rule factors bounds coefficient error by
\[
B_{\rm int}\left[\frac2{e^{a_tN_t}-1}
 +\frac2{e^{h_{\rm ang}N_\varphi}-1}
 +\frac{4(d-2)e^{-\sigma N_\theta}}{1-e^{-\sigma}}\right]
 \le\epsilon_c/2.
\]
Run \Cref{cp:init:source-routine} at $\varepsilon_v=\delta_{\rm node}$.
Base Euclidean error is below $\delta_{\rm node}/8$, hence image
coordinate error is at most $\delta_{\rm node}$.
Spatial weight mass is at most $A_d$, so each coefficient perturbation
is at most $2A_d\mathcal Y\delta_{\rm node}=\epsilon_c/2$.
Adding the analytic tail gives
$\eta/16+N\mathcal Y\epsilon_c=\eta/8$ for sources and images separately.
Accumulate only base coefficients and apply initialized matrices afterward:
the identical scalar sums ensure
$\widetilde c^{\,W_0g}=W_0\widetilde c^{\,g}$ exactly.
Thus no coordinate-error propagation through a matrix norm is needed;
\Cref{cp:selected} applies with $R=2m+d+1+4N$.

For $d=1$, use the two points, $N_x=H_{\rm sph}=2$,
$H_1=\log(64M_n/(\alpha\eta))$, $N_1=1+\lfloor H_1/\alpha\rfloor$,
$p_t=N_1-1$, $N=2N_1$, and $R=2m+d+1+8N_1$.
Take $\epsilon_c=\eta/(16N_1)$,
$\delta_{\rm node}=\eta/(64N_1)$, $B_{\rm int}=2M_ne^{a_tp_t}$,
and $N_t=\max\{1,\lceil a_t^{-1}\log(1+4B_{\rm int}/\epsilon_c)\rceil\}$.
Only the time trapezoid is used. Its error allocation is identical;
the one-dimensional analytic tail is at most $\eta/8$ by
\Cref{cp:harmonic-approximation}, still leaving total error below
$\eta$.

The count proof of \Cref{cp:harmonic} applies to this identical
mode set and the same logarithmic tolerance. In particular its factor
$\alpha^{-1}=128\beta^{30L}(Y/\lambda)^2\sqrt{d+3}\ell^{3/2}$
is preserved before squaring the rank. This proves the complete retained
bound \eqref{cp:harmonic-size}, with its label and dimension factors.

\subsection{Exact implicit Gaussian execution}
\label{cp:init:sec:implicit}

The next finite-dimensional sampler is useful because the numerical
construction accesses an initialized hidden matrix only through its actions.
All its zero tests and normalizations are exact.

\begin{lemma}[Adaptive queries to both directions of a Gaussian matrix]
\label{cp:init:gaussian-sampler}

For each of finitely many layers let the reference matrix have independent
$N(0,1/n)$ entries, independently across layers and of the client's
independent randomness. The client may choose the next layer, direction and
query vector measurably from that randomness and its complete past
transcript, and may stop within deterministic finite query budgets.
There is an online sampler whose entire transcript has exactly the dense
reference law. On an extension of its probability space, independent
Gaussian reference matrices can be completed so that every returned reply
is simultaneously the requested product with that one matrix.

For $Q$ requests to one layer, the sampler uses at most
$4nQ^2+6nQ+2$ scalar arithmetic operations, at most
$n\min\{Q,2n-1\}$ standard normal samples, and
$2n\min\{Q,2n\}+8n+5\min\{Q,2n\}+64$ real or address words.
These counts exclude the client's own work and retained arrays. The
completion matrices are proof objects, not arrays constructed by the algorithm.
\end{lemma}

\begin{proof}
For one layer store $V,Y\in\R^{n\times p}$ and $U,R\in\R^{n\times q}$,
initially empty, with
\[
V^\top V=I,\qquad U^\top U=I,\qquad U^\top Y=R^\top V,
\]
representing $Y=WV$, $R=W^\top U$. For a forward query $v$, compute
$a=V^\top v$, $v_\perp=v-Va$, $\alpha=\|v_\perp\|_2$.
If $\alpha=0$ return $Ya$; otherwise set $e=v_\perp/\alpha$ and
\begin{equation}
\label{cp:init:eq:sampler-forward}
h=UR^\top e+n^{-1/2}[g-U(U^\top g)],
\end{equation}
return $Ya+\alpha h$ and append $(e,h)$ to $(V,Y)$.
Here $g$ is a fresh standard Gaussian vector if $q<n$; if $q=n$
omit the Gaussian term. Transpose queries use this identical algorithm
with $(V,Y,p)\leftrightarrow(U,R,q)$, hence innovation
\begin{equation}
\label{cp:init:eq:sampler-transpose}
h=VY^\top f+n^{-1/2}[g-V(V^\top g)]
\end{equation}
for a normalized new direction $f$, and no draw if $p=n$.
Each layer has independent draws and its own four column lists; no
dense projector is formed.

For proof only put $P=VV^\top$, $Q_0=UU^\top$ and
$A=YV^\top+UR^\top(I-P)$. Compatibility implies $AV=Y$, $A^\top U=R$.
The homogeneous observation space is exactly
$\{(I-Q_0)B(I-P):B\in\R^{n\times n}\}$, orthogonal to $A$.
Inductively the conditional law given the full transcript is
\begin{equation}
\label{cp:init:eq:sampler-kernel}
W=A+(I-Q_0)G(I-P),\qquad G_{ij}\ {\rm independent}\ N(0,1/n).
\end{equation}
It holds initially. Given the past, the adaptive query is fixed.
For a new forward direction set $P'=P+ee^\top$ and split
\[
(I-Q_0)G(I-P)=((I-Q_0)Ge)e^\top+(I-Q_0)G(I-P').
\]
The two terms are independent: their cross-covariances are
$n^{-1}(I-Q_0)_{ia}[e^\top(I-P')]_b=0$.
This remains true for singular Gaussians, since the joint characteristic
function of linear images factors when their cross-covariance vanishes.
Thus $We$ has the sampler law, and conditioning updates the mean to
$A+(h-UR^\top e)e^\top$, precisely the new stored-state expression,
while $U^\top h=R^\top e$ preserves compatibility.
Transposition exchanges $P,Q_0$ and $(V,Y),(U,R)$ in this calculation:
with $Q_0'=Q_0+ff^\top$, the independent split is
$f((I-P)G^\top f)^\top+(I-Q_0')G(I-P)$ and the mean update is
$A+f(h-VY^\top f)^\top$.
Dependent queries leave the kernel unchanged; full opposite rank removes
the innovation.

Conditional layer kernels remain a product: a transcript-measurable
choice of layer or direction reveals nothing further, and its reply
conditions only that layer. Finite induction, partitioned by the finitely
many stopping events if needed, proves equality of transcript laws.
At stopping, use fresh independent $G$ in
\eqref{cp:init:eq:sampler-kernel} for each layer. Integrating these
conditional kernels gives the ordinary joint Gaussian matrix law,
including independence from client randomness. The completion agrees on
all final basis columns, hence on every earlier reply, pathwise.

At ranks $p,q$, direct dot-product counts give at most
$6np+8nq+9n+1\le8n(p+q)+10n$ operations per forward request;
transpose swaps the ranks and dependent queries cost less.
Since $p+q\le j-1$ before request $j$, summation gives
$4nQ^2+6nQ$, plus two operations for $1/\sqrt n$.
There are at most $2n-1$ drawing increments before a rank reaches $n$;
afterward every new opposite action is determined.
The four lists occupy $2n(p+q)$ words, with $p+q\le\min(Q,2n)$.
Eight length-$n$ scratch vectors, five coefficient/address words per
column and 64 scalar words give the stated memory. At most two
comparisons per request suffice. Exact zero tests require neither a
condition-number bound nor a nonsingular query Gram.
\end{proof}

\subsubsection{Factored numerical anchors and output coupling}
\label{cp:init:sec:factored-execution}

Fix all orders, scalar routines, generator ordering, orthogonalization
signs and selector tie-breaking. On panel $b$ of length $\Delta_b$,
use unscaled coefficients in $t-\tau_b$. In the polynomial network put
$U_i^{(j)}=([r_a\delta_a^{(j)}]_i)_{a=1}^m$ and
$H_k^{(j)}=([h_a^{(j)}]_k)_{a=1}^m$, both $n$ by $m$ matrices.
Coefficient comparison in the actual gradient equations gives, for $s\ge1$,

\begin{equation}
\label{cp:init:eq:factored-jets}
\begin{aligned}
[W^{(1)}]_s&=-\frac2{ms}U_{s-1}^{(1)}(v_a)_{a=1}^m{}^\top,\\
[W^{(j)}]_s&=-\frac2{mns}\sum_{i+k=s-1}U_i^{(j)}H_k^{(j-1)\top},\\
[w]_s&=-\frac2{ms}\sum_a[r_ah_a^{(L)}]_{s-1}.
\end{aligned}
\end{equation}

Only lower orders determine order $s$. Acting on a query series $X$ gives

\begin{equation}
\label{cp:init:eq:cached-action}
-\frac2{mn}\sum_{i+k+c=s-1}
             \frac{U_i^{(j)}(H_k^{(j-1)\top}X_c)}{i+k+1}.
\end{equation}

Cache $H_k^\top X_c$: for $b'$ queries, vector work is
$O(nmb'K^2)$ and scalar sums cost $O(mb'K^3)$ per layer.
Transposition exchanges the factors. Training coefficients are causal;
passive jets use the completed training jets.

The endpoint update groups into $mK$ outer products per hidden layer:

\begin{equation}
\label{cp:init:eq:anchor-increment}
W_{b+1}^{(j)}-W_b^{(j)}
=-\frac2{mn}\sum_{i=0}^{K-1}U_i^{(j)}V_i^{(j-1)\top},\qquad
V_i^{(j-1)}=\sum_{k=0}^{K-1-i}
             \frac{\Delta_b^{i+k+1}}{i+k+1}H_k^{(j-1)}.
\end{equation}

Explicit execution adds these outer products to the dense anchor;
implicit execution appends them, with the scalar in $A_b$, to
$A_b^{(j)},B_b^{(j)}\in\R^{n\times bmK}$:

\begin{equation}
\label{cp:init:eq:implicit-anchor}
W_b^{(j)}=W_0^{(j)}+A_b^{(j)}B_b^{(j)\top},\qquad
W_b^{(j)}X_c=\operatorname{Action}_j(X_c)
                  +A_b^{(j)}(B_b^{(j)\top}X_c).
\end{equation}

Transpose actions exchange the factors. This exactly represents numerical
anchors, without asserting low rank of the continuous dense trajectory.
The first matrix is explicit and its higher coefficients use training
input pairings in \eqref{cp:init:eq:factored-jets}, without $K$ dense
first-matrix arrays; the readout is a vector. Hidden initialized-matrix
requests are exactly:

\begin{enumerate}
\item $m$ forward actions per mixer for exact initialized training
features/images using $\phi_j$; initial backward fields vanish.
\item Constant-anchor actions on source/training coefficients through $K$
per panel; all increments use the displayed factors.
\item Images of completed Harmonic coefficients or Taylor nodes.
\item Images $W_0^{(j)}U_{j-1}$ of the final source basis, at most
$r$ per mixer, including any not spanned by earlier image requests.
\end{enumerate}
Harmonic uses both directions on $m+N_x$ inputs; Taylor uses $M$
forward and $m$ backward inputs. Sufficient per-mixer counts are

\begin{equation}
\label{cp:init:eq:query-counts}
Q_{\rm H}=m+2J(m+N_x)(K+1)+2N+r,\qquad
Q_{\rm T}=m+J(M+m)(K+1)+(M+m)N_s+r.
\end{equation}

Use deterministic $\min(n,R)$ in place of $r$ to invoke the sampler;
the realized cost may use actual $r$. No other hidden-matrix information
(entries, norms or singular values) is accessed.
Under completion, explicit and implicit executions agree by induction
through layers, coefficients and panels, using
\cref{cp:init:eq:implicit-anchor,cp:init:eq:cached-action}.
Identical generator lists then give identical ranks, bases and selector
decisions. For source bases
$U_j^\top U_j/n=I$, selected row matrices $P_j=(U_j)_{I_j}$,
and source metrics $M_j$, item 4 gives the same core matrix
$U_j^\top W_0^{(j)}U_{j-1}/n$ and thus the same initialized mixer

\[
B_C^{(j)}(0)=P_j\frac{U_j^\top W_0^{(j)}U_{j-1}}n
                         P_{j-1}^\top M_{j-1}.
\]

The first weights, zero raw readout, $c_C=y$, exact initial Gram and
corrected-readout caches also agree. Hence retained arrays and autonomous
futures agree pathwise, with the same source/fitting probability and no
extra failure allowance. Transcripts and trained factors are discarded.

\subsection{Complete operation and storage ledgers}
\label{cp:init:sec:costs}

Finite-order big-$O$ constants below are numerical: no problem parameter,
order, scalar-evaluator cost or confidence is hidden. Assume
$n\ge\max(m,d)$; full-width retention is available otherwise when
the source budget reaches $n$. Let $R$ bound generators, $r$ be maximum
actual rank and $q$ maximum selected support, so
$r\le\min(n,R)$ and $r\le q\le9r$. These are local cost
quantities, distinct from the source-dimension argument of \Cref{cp:selected}.

\subsubsection{Scalar geometry, global coefficients, and selection}
\label{cp:init:sec:geometry-cost}

For Harmonic, assign each time node $t_a=T(1+\cos u_a)/2$ to
one numerical panel, retaining multiplicity if two time angles have the
same physical time. Let $I_b$ be those node indices. Form

\[
\mathcal B_{kc}^{(b)}=\frac{\gamma_k}{N_t}\sum_{a\in I_b}
 \cos(ku_a)((t_a-\tau_b)/\Delta_b)^c.
\]

If $[g]_c(v_s)$ is an unscaled local source coefficient, first multiply
it by $\Delta_b^c$ to obtain $G_{b,c}(v_s)$. Its global coefficient
is exactly

\begin{equation}
\label{cp:init:eq:streamed-projection}
\widetilde c_{kjb'}=
 \sum_b\sum_s\omega_sY_{jb'}(v_s)
                  \sum_{c=0}^K\mathcal B_{kc}^{(b)}G_{b,c}(v_s).
\end{equation}

Process each panel's training jets, spatial queries and anchor advance
in that order; accumulate spatial coefficients before transforming retained
mode pairs. Scalar tables cost
$O((N_t+J)(p_t+1)(K+1))$ work, and projection costs
$O(LnJ(K+1)(N_xH_{\rm sph}+N))$ work with
$O(Ln(K+1)H_{\rm sph})$ extra vector words. Only $O(LnR)$
global coefficient words persist across panels. Local panels do not
create separate retained Harmonic source blocks.

For finite basis evaluation, use the circle basis
$1,\sqrt2\cos(j\varphi),\sqrt2\sin(j\varphi)$ and angle addition.
For $d'\ge3$, at $v=(\cos\theta,\sin\theta\,u)$ use normalized

\[
(\sin\theta)^l C_{j-l}^{\,l+(d'-2)/2}(\cos\theta)
                       Y_{lb}(u),\qquad 0\le l\le j,
\]

Sphere-Laplacian substitution gives eigenvalue $-j(j+d'-2)$.
For each orthogonal lower harmonic, the remaining weight is
$(1-z^2)^{l+(d'-3)/2}$. Multiplying the Gegenbauer equation
$(1-z^2)y''-(2\mu+1)zy'+s(s+2\mu)y=0$ by
$(1-z^2)^{\mu-1/2}$ and integrating by parts proves orthogonality:
$\mu\ge1/2$ kills the boundary term. The dimension decomposition
in \Cref{cp:harmonic-approximation} proves completeness. Use

\[
(s+1)C_{s+1}^{\mu}(z)=2(s+\mu)zC_s^\mu(z)
 -(s+2\mu-1)C_{s-1}^\mu(z),\quad C_0^\mu=1,\quad C_1^\mu=2\mu z.
\]

from the generating function $(1-2zt+t^2)^{-\mu}$.
For squared angular norm $N_{l,s}$ under probability-sphere measure,
$\mu=l+(d'-2)/2$ gives

\[
N_{0,0}=1,\quad
\frac{N_{l+1,0}}{N_{l,0}}=\frac{2l+d'-1}{2l+d'},\qquad
\frac{N_{l,s+1}}{N_{l,s}}=
 \frac{(s+2\mu)(s+\mu)}{(s+1)(s+\mu+1)}.
\]

The first ratio follows by integration by parts on sine powers; the
second follows by pairing the recurrence for $zC_s^\mu$ with
$C_{s+1}^\mu$ and using symmetry of multiplication by $z$.
Only recurrences and inverse square roots are required. Per intermediate
dimension there are $O((j_*+1)^2)$ degree pairs and $O(H_{\rm sph})$
lower-basis products. Recompute spatial values per panel, giving

\[
\mathcal G_{\rm time}=\begin{cases}
N_t+J,&d=1,\\
N_t+N_x+JN_x(j_*+1),&d=2,\\
N_\theta^2+dN_\theta+N_t+N_\varphi+d+d(j_*+1)^2
 +JdN_x[1+(j_*+1)^2+H_{\rm sph}],&d\ge3,
\end{cases}
\]

\begin{equation}
\label{cp:init:eq:geometry-counts}
\mathcal G_{\rm memory}=\begin{cases}
N_t+1,&d=1,\\
N_t+N_x+j_*+1,&d=2,\\
dN_\theta+N_t+N_\varphi+d[1+(j_*+1)^2+H_{\rm sph}],&d\ge3.
\end{cases}
\end{equation}

These bounds include $O(N_\theta^2)$ for \eqref{cp:init:eq:fejer},
$O(dN_\theta)$ for sine powers, and $O(d)$ per node for coordinates
and Jacobians, as well as the displayed basis recurrence.

For either model, orthogonalizing $R$ generators costs $O(LnRr)$.
Each of the at most $9r$ selector steps from \Cref{cp:selection}
costs $O(r^3+nr^2)$ for shifted inverses and all-row quadratic forms.
Thus selection costs $O(Lnr^3)$ since $r\le n$.

For the final metric, suppress layer indices and let $P$ be the
selected basis rows, $\mathsf D$ the positive diagonal selector
weights and $G=P^\top\mathsf DP$. Formula \eqref{cp:metric}
and direct multiplication give

\[
M=\mathsf D+\mathsf DP(G^{-2}-G^{-1})P^\top\mathsf D,\quad
M^{-1}=\mathsf D^{-1}+P(I-G^{-1})P^\top,\quad
P^\top M=G^{-1}P^\top\mathsf D.
\]

The $r$ by $r$ Gram has spectrum in $[1,4]$. Metric caches and
core mixers cost $O(L(qr^2+r^3+q^2r))=O(Lq^2r)$.
Basis images cost $O(n^2r)$ per explicit mixer, or requests in $Q$
plus $O(nr^2)$ for implicit contraction. The initial training Gram,
identical by exact source isometry, costs $O(nm^2+m^3)$ work and
$O(m^2)$ words including factorization.

Degree-$D$ activation composition with degree-$K$ series costs
$O(DK^2)$ work and $O(DK)$ scratch: the $O(D)$ Chebyshev
products each cost $O(k+1)$ at coefficient $k$.
Preparing coefficients and derivatives costs $O(LD^2)$ once.
Training and passive sources use this same routine.

\subsubsection{Shared finite execution costs}
\label{cp:init:sec:explicit-costs}
\label{cp:init:sec:implicit-costs}

The four ledgers share the same continuation, orthogonalization and
selection routines. To expose that reuse without suppressing parameters,
let $e$ be the explicit input dimension and $B,F$ local query-batch
sizes, and put $P_e=(L-1)n^2+n(e+1)$. Define the exact sums
\begin{align*}
\mathcal W_{\rm ex}(P,B)={}&LD^2+mP+nm^2+m^3
 +J\{PBK+LmBK^2(n+K)+LnBDK^2\}\\
 &+LnRr+Lnr^3+Lq^2r,\\
\mathcal M_{\rm ex}={}&P_d+LmnDK+LnR+Lq^2+LD+J+m^2+L,\\
\mathcal W_{\rm im}(e,F,Q)={}&ne(1+m)+em^2+nm^2+m^3+L+e+m+1+L(D+1)^2\\
 &+(L-1)nQ^2+(L-1)nm(m+F)J(J-1)K(K+1)\\
 &+J[ne(m+F)+emF]+JLm(m+F)(K+1)^2(n+K+1)\\
 &+JLn(m+F)(D+1)(K+1)^2+LnRr+Lnr^3\\
 &+(L-1)nr^2+L(qr^2+r^3+q^2r)+qe,\\
\mathcal M_{\rm im}(e,Q)={}&ne+(L-1)n[Q+JmK]+Lmn(D+1)(K+1)\\
 &+LnR+Lq^2+q(e+Lm+1)+m^2+L(D+1)+J+L+e+m+1.
\end{align*}
The substitutions are $(e,F)=(d,N_x)$ for Harmonic and $(k,M)$
for Taylor; in the latter $M$ includes training and the extra $m$
pays backward queries. These symbols are local cost arguments.
In each display below the work or peak memory is bounded by $O$ of the
entire right-hand sum, with numerical constants only.

For supplied matrices the Harmonic and Taylor work bounds are
\begin{align}
&\mathcal W_{\rm ex}(P_d,m+N_x)+P_d+L+d+m+1
 +(N_t+J)(p_t+1)K \nonumber\\
&\qquad+LnJK(N_xH_{\rm sph}+N)+Ln^2(N+r)
 +\mathcal G_{\rm time}+N_t\log(2+N_t),
\label{cp:init:eq:explicit-harmonic-work}\\
&\mathcal W_{\rm ex}(P_k,M+m)+P_d+dM^2+ndk+L+k+m+1 \nonumber\\
&\qquad+LnMKN_s+Ln^2R+N_s\log(2+N_s)+J_s+N_s,
\label{cp:init:eq:explicit-taylor-work}
\end{align}
and their respective peak-memory sums are
\begin{align}
&\mathcal M_{\rm ex}+Lm^2K^2+LnKH_{\rm sph}
 +(p_t+1)K+N_t+m(d+1)+d+\mathcal G_{\rm memory},
\label{cp:init:eq:explicit-harmonic-memory}\\
&\mathcal M_{\rm ex}+LnMK+LmMK^2+N_s+J_s+Md+dk+k+m+1.
\label{cp:init:eq:explicit-taylor-memory}
\end{align}
Constant-anchor products cost $JPBK$; cached higher coefficients and
activation compositions are the next two terms of $\mathcal W_{\rm ex}$.
Grouping \eqref{cp:init:eq:anchor-increment} costs $O(nmK^2)$ and
adding its $mK$ outer products costs $O(n^2mK)$ per layer and panel,
already covered. Harmonic separately pays $Ln^2N$ for source images
and $Ln^2r$ for basis images; Taylor bounds both by $Ln^2R$.
The Taylor reduction costs $dM^2+ndk$. Sort source nodes into numerical
panels and evaluate them in one pass for $O(LnMKN_s)$.
One anchor, original matrices, current jets and global coefficients are
live; neither a $JP_d$ trajectory nor $KP_d$ dense jets are stored.

For fresh implicit matrices set $Q=Q_{\rm H}$ or $Q_{\rm T}$ from
\eqref{cp:init:eq:query-counts}. The corresponding work sums are
\begin{align}
&\mathcal W_{\rm im}(d,N_x,Q_{\rm H})+(N_t+J)(p_t+1)(K+1)\nonumber\\
&\qquad+LnJ(K+1)(N_xH_{\rm sph}+N)
 +\mathcal G_{\rm time}+N_t\log(2+N_t),
\label{cp:init:eq:implicit-harmonic-work}\\
&\mathcal W_{\rm im}(k,M,Q_{\rm T})+dM^2
 +Ln(M+m)(K+1)N_s\nonumber\\
&\qquad+N_s\log(2+N_s)+N_s+J+J_s,
\label{cp:init:eq:implicit-taylor-work}
\end{align}
and their peak-memory sums are
\begin{align}
&\mathcal M_{\rm im}(d,Q_{\rm H})+Lm^2(K+1)^2+Ln(K+1)H_{\rm sph}\nonumber\\
&\qquad+m(d+1)+(p_t+1)(K+1)+N_t+\mathcal G_{\rm memory},
\label{cp:init:eq:implicit-harmonic-memory}\\
&\mathcal M_{\rm im}(k,Q_{\rm T})+Md+dk+LnM(K+1)
 +LmM(K+1)^2+Mk+N_s+J_s.
\label{cp:init:eq:implicit-taylor-memory}
\end{align}
The sampler contributes $(L-1)nQ^2$ work and $(L-1)nQ$ memory.
At panel $b$, the $bmK$ trained-factor columns cost $O(nbmK)$
per coefficient action. Summing at most $2(m+F)(K+1)$ actions
and $\sum_{b<J}b=J(J-1)/2$ gives the second work term;
the factor memory is $(L-1)nJmK$. Their grouped formation is
covered by current-panel contractions. All image queries, the $nr^3$
selector, and final Gram and metric solves are charged explicitly.

Implicit Taylor samples only the $n$ by $k$ reduced first matrix.
Its law is $W_0^{(1)}U_0$; an independent Gaussian component on the
orthogonal input complement is completed only in the proof and is never
queried. These finite sums therefore charge every accessed matrix action
and every trained-factor product, independently of the eventual rank.

\subsubsection{Primitive costs and fixed-parameter bounds}
\label{cp:init:sec:specialization}

Each model uses $4L(D+1)+Lnm$ original activation-value calls,
plus at most $L$ unsupplied values at zero. Only the first $4L(D+1)$
permit the tolerance of \eqref{cp:init:eq:activation-polynomial}.
Setup needs no original derivatives; runtime uses original values and
first derivatives.

Supplied matrices need no draws; fresh explicit initialization adds
$nd+(L-1)n^2$. Implicit Harmonic and Taylor use at most
$nd+(L-1)nQ_{\rm H}$ and $nk+(L-1)nQ_{\rm T}$ respectively,
with no readout or proof-completion draws. Harmonic's elementary-function
call count is at most

\[
O(L+d+m+1+L(D+1)+N_t+A_{\rm ang}+(L-1)Q+Lr),\quad
A_{\rm ang}=\begin{cases}
0,&d=1,\\N_x,&d=2,\\
N_\theta+N_\varphi+d(j_*+1)^2,&d\ge3.
\end{cases}
\]

Taylor replaces time/angular terms by $N_s+J_s$ and adds $O(M)$
input-basis square roots; explicit execution omits $Q$.
These counts cover certificates, orders, geometry and normalizations;
rounding/enumeration fit the finite tables. Nonprimitive scalar calls add
their actual work, description and peak scratch. Envelopes and population
gap certificates are supplied, not computed from analyticity.

For separately fixed admissible parameters,
\eqref{cp:init:eq:numerical-scales} and the source formulas give

\begin{equation}
\label{cp:init:eq:all-orders}
\begin{array}{c|c|c}
&\text{Harmonic}&\text{Taylor}\\ \hline
J,D& O(\ell^{3/2})&O(\ell^{3/2})\\
K&O(\ell)&O(\ell)\\
\text{source sizes}&N_t,p_t+1=O(\ell^{5/2}),\quad
 N_x,H_{\rm sph}=O(\ell^{3(d-1)/2})&
 J_s=O(\ell^{1/2}),\quad K_s+1=O(\ell)\\
R,r,q&O(\ell^{3d/2+1})&O(\ell^{3/2})\\
Q&O(\ell^{3d/2+1})&O(\ell^{5/2}).
\end{array}
\end{equation}

The two-point convention gives the same Harmonic exponents at $d=1$.
For the supplied dense variants the full ledgers respectively imply

\[
O(n^2\ell^{3d/2+1}+n\ell^{9d/2+3}+\operatorname{polylog}(n)),\qquad
O(n^2\ell^{5/2}+n\ell^5+\operatorname{polylog}(n)).
\]

At fixed parameters the leading terms dominate eventually, and memory
is $O(n^2)$. For implicit Harmonic, sampler, trained-factor/activation
and projection exponents are respectively $3d+2$, $3d/2+7/2$,
and at most $\max(3d-1/2,3d/2+7/2)$. For $d\ge1$ all are
dominated by the selector $O(nr^3)=O(n\ell^{9d/2+3})$;
non-width terms are fixed powers of $\ell$ and memory is
$O(n\ell^{3d/2+1})$. For implicit Taylor,

\[
nQ_{\rm T}^2=O(n\ell^5),\qquad
nm(M+m)J^2K^2=O(n\ell^5),\qquad
Jn(M+m)DK^2=O(n\ell^5).
\]

The selector costs $O(n\ell^{9/2})$ and all remaining work fits
$O(n\ell^5)$, while transcript, factors, compositions and other
buffers fit $O(n\ell^{5/2})$. This proves all four rows of
\Cref{cp:initialization} including final assembly.

Discard anchors, generators, bases, schedules, activation polynomials,
transcripts and factors; retain exactly the inventory and optimizer of
\Cref{cp:selected}. Selection uses the computed span and bound $R$,
so accidental rank deficiencies need not survive approximation.
Exact initialized additions preserve the Gram and pairing. Finite-precision
values, ranks, barrier comparisons or infinite-time evolution would need
new perturbation results: these are real-coordinate and operation counts.
The setup covers $T=32\ell/\lambda$, not merely a short physical prefix.
Dense Euler at step $h$ costs $\Theta(mP_dT/h)$ before assembly;
the fixed-parameter Taylor/Harmonic setup ratios are
$O(h\ell^{3/2})$, $O(h\ell^{3d/2})$. This neither proves that $h$
is necessary nor compares against every high-order dense solver.
